\documentclass[11pt,a4paper]{article}

% ── Codificación y fuentes ────────────────────────────────────────────────────
\usepackage[utf8]{inputenc}
\usepackage[T1]{fontenc}
\usepackage{lmodern}
\usepackage[spanish]{babel}

% ── Geometría y espaciado ─────────────────────────────────────────────────────
\usepackage[top=2.5cm, bottom=2.5cm, left=2.8cm, right=2.8cm]{geometry}
\usepackage{setspace}
\setstretch{1.15}
\usepackage{parskip}

% ── Tablas ────────────────────────────────────────────────────────────────────
\usepackage{booktabs}
\usepackage{tabularx}
\usepackage{array}
\usepackage{longtable}
\usepackage{enumitem}

% ── Colores y cajas ───────────────────────────────────────────────────────────
\usepackage[dvipsnames]{xcolor}
\usepackage{tcolorbox}
\tcbuselibrary{skins, breakable}

% ── Cabeceras y pies de página ────────────────────────────────────────────────
\usepackage{fancyhdr}
\usepackage{lastpage}
\pagestyle{fancy}
\fancyhf{}
\fancyhead[L]{\small\textbf{Informe de Evaluación} --- Física para Bioanalistas}
\fancyhead[R]{\small yilselys betancourt 33385584}
\fancyfoot[C]{\small Página \thepage\ de \pageref{LastPage}}
\renewcommand{\headrulewidth}{0.4pt}

% ── Hiperenlaces ──────────────────────────────────────────────────────────────
\usepackage[colorlinks=true, linkcolor=NavyBlue, urlcolor=NavyBlue]{hyperref}

% ── Paleta de colores personalizados ─────────────────────────────────────────
\definecolor{desempeno_excelente}{HTML}{1A6B2F}
\definecolor{desempeno_bueno}{HTML}{2E5A9C}
\definecolor{desempeno_suficiente}{HTML}{8B6914}
\definecolor{desempeno_deficiente}{HTML}{C0392B}
\definecolor{desempeno_insuficiente}{HTML}{7B241C}
\definecolor{grisFondo}{HTML}{F5F5F5}
\definecolor{azulTitulo}{HTML}{1B2A6B}
\definecolor{verdeFortaleza}{HTML}{1A6B2F}
\definecolor{naranjaMejora}{HTML}{A04000}

% ── Estilos de cajas ─────────────────────────────────────────────────────────
\tcbset{
  cajaTitulo/.style={
    enhanced, breakable,
    colback=azulTitulo!8, colframe=azulTitulo,
    fonttitle=\bfseries\large, coltitle=white,
    attach boxed title to top left={yshift=-2mm, xshift=4mm},
    boxed title style={colback=azulTitulo},
    top=6pt, bottom=6pt, left=8pt, right=8pt,
  },
  cajaFortaleza/.style={
    enhanced, breakable,
    colback=verdeFortaleza!6, colframe=verdeFortaleza!60,
    fonttitle=\bfseries, coltitle=verdeFortaleza,
    top=4pt, bottom=4pt, left=8pt, right=8pt,
  },
  cajaMejora/.style={
    enhanced, breakable,
    colback=naranjaMejora!6, colframe=naranjaMejora!60,
    fonttitle=\bfseries, coltitle=naranjaMejora,
    top=4pt, bottom=4pt, left=8pt, right=8pt,
  },
  cajaRecomendacion/.style={
    enhanced, breakable,
    colback=grisFondo, colframe=gray!50,
    fonttitle=\bfseries, coltitle=black,
    top=4pt, bottom=4pt, left=8pt, right=8pt,
  },
}

% ─────────────────────────────────────────────────────────────────────────────
\begin{document}

% ══ PORTADA ══════════════════════════════════════════════════════════════════
\begin{center}
  {\Large\bfseries\color{azulTitulo} Universidad de Oriente/Núcleo Sucre --- Física para Bioanalistas}\\[4pt]
  {\large Informe de Evaluación Preliminar}\\[12pt]
  \rule{\linewidth}{1.2pt}\\[6pt]
  {\LARGE\bfseries yilselys betancourt 33385584}\\[4pt]
  \rule{\linewidth}{0.6pt}
\end{center}

% ══ FICHA TÉCNICA ════════════════════════════════════════════════════════════
\vspace{4pt}
\begin{tcolorbox}[cajaTitulo, title=Datos del informe]
\begin{tabular}{@{}llll@{}}
  \textbf{Estudiante:}  & yilselys betancourt 33385584  &
  \textbf{Fecha:}       & 2026-06-26 \\[4pt]
  \textbf{Archivo PDF:} & \texttt{yilselys\_betancourt\_33385584.pdf}  &
  \textbf{Páginas:}     & 4 \\
\end{tabular}
\end{tcolorbox}

% ══ RESULTADO GLOBAL ═════════════════════════════════════════════════════════
\vspace{6pt}
\begin{center}
  \begin{tcolorbox}[
    enhanced,
    colback=white, colframe=desempeno_bueno,
    width=0.62\linewidth,
    boxrule=2pt,
    halign=center, valign=center,
    top=8pt, bottom=8pt,
  ]
    \centering
    {\large\bfseries Puntaje sugerido}\\[6pt]
    {\Huge\bfseries\color{desempeno_bueno} 8.3/10}\\[6pt]
    {\large Nivel de desempeño: \textbf{\color{desempeno_bueno} Bueno}}
  \end{tcolorbox}
\end{center}

% ══ RESUMEN GENERAL ══════════════════════════════════════════════════════════
\section*{Resumen general}
El estudiante demuestra un buen entendimiento de los principios físicos fundamentales aplicados en el examen, especialmente en hidrodinámica (ecuación de continuidad, Bernoulli y Poiseuille). La identificación de datos y la selección de fórmulas son generalmente correctas. Sin embargo, se observan algunas imprecisiones en la identificación inicial de variables y un error significativo en el cálculo de exponentes en la última pregunta, lo que afectó el resultado final. La organización y claridad de las respuestas son adecuadas.

% ══ FORTALEZAS Y ÁREAS DE MEJORA ════════════════════════════════════════════
\vspace{4pt}
\begin{minipage}[t]{0.48\linewidth}
  \begin{tcolorbox}[cajaFortaleza, title={Fortalezas}]
\begin{itemize}[leftmargin=1.5em, itemsep=2pt]
  \item Correcta aplicación de la ecuación de continuidad y Bernoulli para el cálculo de velocidad y presión.
  \item Excelente comprensión y aplicación de la Ley de Poiseuille para análisis proporcional.
  \item Buena capacidad para identificar los datos relevantes en la mayoría de los problemas y realizar conversiones de unidades cuando es necesario.
  \item Organización clara de los pasos de resolución (Datos, Fórmula, Solución, Respuesta).
  \item Manejo adecuado de unidades en la mayoría de los cálculos.
\end{itemize}
  \end{tcolorbox}
\end{minipage}
\hfill
\begin{minipage}[t]{0.48\linewidth}
  \begin{tcolorbox}[cajaMejora, title={Areas de mejora}]
\begin{itemize}[leftmargin=1.5em, itemsep=2pt]
  \item Precisión en la identificación y etiquetado de variables al inicio de los problemas para evitar confusiones.
  \item Revisión cuidadosa de los cálculos, especialmente con exponentes y notación científica.
  \item Asegurarse de responder completamente a la pregunta formulada, incluyendo todos los pasos requeridos para el resultado final.
\end{itemize}
  \end{tcolorbox}
\end{minipage}

% ══ OBSERVACIONES POR PREGUNTA ═══════════════════════════════════════════════
\section*{Observaciones por pregunta}

\begin{longtable}{>{\bfseries}p{2.8cm} p{1.6cm} p{7.0cm} p{2.8cm}}
  \toprule
  \textbf{Pregunta} & \textbf{Puntaje} & \textbf{Evaluación} & \textbf{Errores detectados} \\
  \midrule
  \endhead
  \bottomrule
  \endfoot
    \textbf{Pregunta 1: Presión Hidrostática y Venosa} & 1.8/2.0 & El estudiante identificó correctamente la fórmula para la presión hidrostática y la despejó adecuadamente para encontrar la altura. Los valores numéricos utilizados en el cálculo son correctos y el resultado final es preciso. Sin embargo, en la sección de 'Datos', hubo una confusión al etiquetar las variables 'h' y 'P', asignando el valor de la presión a 'h' y viceversa, aunque esto no afectó el uso correcto de los valores en la fórmula. & \textit{Confusión en el etiquetado de variables (h y P) en la sección de datos.} \\
    \midrule
    \textbf{Pregunta 2: Principio de Arquímedes} & 1.5/2.0 & El estudiante aplicó correctamente el principio de Arquímedes para calcular la fuerza de empuje y, a partir de ella, el volumen del tejido óseo. Los cálculos para la fuerza de empuje y el volumen son correctos. Sin embargo, la pregunta solicitaba la densidad del tejido óseo, y el estudiante se detuvo en el cálculo del volumen, dejando la respuesta incompleta. & \textit{Respuesta incompleta: no se calculó la densidad final del tejido óseo, solo su volumen.} \\
    \midrule
    \textbf{Pregunta 3: Hemodinámica (Continuidad y Bernoulli)} & 2.0/2.0 & El estudiante resolvió esta pregunta de manera excelente. Aplicó correctamente la ecuación de continuidad para calcular la velocidad en la zona de estenosis y luego utilizó la ecuación de Bernoulli (asumiendo flujo horizontal) para determinar la presión en dicha zona. Todos los cálculos son precisos y las unidades son consistentes. & \textit{---} \\
    \midrule
    \textbf{Pregunta 4: Ley de Poiseuille (Análisis Proporcional)} & 2.0/2.0 & La resolución de esta pregunta es impecable. El estudiante comprendió y aplicó correctamente la relación de proporcionalidad de la Ley de Poiseuille (Q $\propto$ r$^{4}$) para determinar el nuevo caudal en función de la reducción del radio. El cálculo y la expresión del resultado en porcentaje son correctos. & \textit{---} \\
    \midrule
    \textbf{Pregunta 5: Flujo Viscoso y Resistencia} & 1.0/2.0 & El estudiante identificó correctamente la Ley de Poiseuille para la diferencia de presión y realizó las conversiones de unidades necesarias (diámetro a radio). La sustitución de los valores en la fórmula es correcta. Sin embargo, se cometió un error significativo en el cálculo del exponente final durante la división, lo que llevó a un resultado numérico incorrecto en varias órdenes de magnitud. & \textit{Error procedimental en el cálculo del exponente final durante la división (10$^{-}$$^{1}$$^{0}$ / 10$^{-}$$^{1}$$^{4}$ se calculó como 10$^{-}$$^{2}$$^{0}$ en lugar de 10$^{4}$).} \\
    \midrule
\end{longtable}

% ══ ERRORES DETECTADOS ═══════════════════════════════════════════════════════
\section*{Análisis de errores}

\subsection*{Errores conceptuales}
\textit{Ninguno identificado.}

\subsection*{Errores procedimentales}
\begin{itemize}[leftmargin=1.5em, itemsep=2pt]
  \item Confusión en el etiquetado de variables en la sección de datos (Pregunta 1).
  \item Respuesta incompleta al no realizar el cálculo final de la densidad (Pregunta 2).
  \item Error en el cálculo de exponentes en la división final, resultando en un valor numérico incorrecto (Pregunta 5).
\end{itemize}

% ══ RECOMENDACIONES AL ESTUDIANTE ════════════════════════════════════════════
\section*{Retroalimentación para el estudiante}
\begin{tcolorbox}[cajaRecomendacion, title={Dirigido a: yilselys betancourt 33385584}]
Estimado estudiante, has demostrado un sólido conocimiento de los principios de la física de fluidos, aplicando correctamente ecuaciones clave como la de continuidad, Bernoulli y Poiseuille. Tu organización en la presentación de las soluciones es muy buena. Para mejorar, te sugiero prestar más atención a la identificación precisa de todas las variables al inicio de cada problema y asegurarte de que tus respuestas aborden completamente lo que se pregunta. Es crucial revisar con detenimiento los cálculos, especialmente aquellos que involucran notación científica y exponentes, ya que un pequeño error puede llevar a un resultado final incorrecto. Sigue practicando y prestando atención a estos detalles para alcanzar la excelencia.
\end{tcolorbox}

% ══ NOTA PARA EL PROFESOR ════════════════════════════════════════════════════
\section*{Nota para el profesor}
\begin{tcolorbox}[
  enhanced, breakable,
  colback=yellow!8, colframe=yellow!60!black,
  fonttitle=\bfseries, coltitle=black,
  title={Observaciones para revisión manual},
  top=4pt, bottom=4pt, left=8pt, right=8pt,
]
El estudiante muestra un buen dominio conceptual de la materia. Los errores son principalmente de tipo procedimental o de atención al detalle (etiquetado de datos, finalización de la respuesta, cálculo de exponentes). La pregunta 5 es un claro ejemplo de un error de cálculo de exponente que desvirtúa el resultado final, a pesar de que la fórmula y la sustitución inicial fueron correctas. La pregunta 2 se penalizó por no completar el cálculo de la densidad final, aunque el volumen se obtuvo correctamente.
\end{tcolorbox}

% ══ PIE DE INFORME ════════════════════════════════════════════════════════════
\vspace{12pt}
\begin{center}
  \footnotesize\color{gray}
  Informe generado automáticamente por el sistema de evaluación con IA (gemini-2.5-flash) y soporte LaTex. \\
  Este documento es una evaluación \textbf{preliminar} para ser revisado por el profesor antes de ser entregado al 
  estudiante. \\
  Generado el 2026-06-26 \textbullet\ BIOANALISIS Sistema Académico de Evaluación.
  \end{center}

\end{document}
